\section{控制收敛用一个可积包络守住质量}
\label{sec:12-Real-Analysis-Support/04-Limits-and-Integration/02}

上一节留了个悬案。\(f_n(x)=x^n\) 在 \([0,1]\) 上不一致收敛，极限函数在 \(x=1\) 处断开，单调收敛定理也用不上——它是递减的。可积分照样穿过了极限号。既然三条已知的通道都不通，是什么在替我们拦住质量？

把两列函数摆在一起对比就看出来了。对每个 \(x\) 和每个 \(n\) 都有 \(0\le x^n\le1\)，常数 \(1\) 在 \([0,1]\) 上可积，一个与 \(n\) 无关的可积函数从上方压住了整族。换成上一节的尖峰列 \(f_n=n\mathbf1_{(0,1/n)}\)，逐点上确界 \(\sup_n f_n(x)\) 在零点附近与 \(1/x\) 同阶，而 \(1/x\) 在 \((0,1]\) 上积分发散。整族头顶没有一片可积的天花板，面积就能一路挤进原点。

\begin{center}
\begin{tikzpicture}[>=stealth]
\draw[->] (0,0) -- (5.9,0) node[right] {\footnotesize \(x\)};
\draw[->] (0,0) -- (0,3.5);
\draw[very thick,domain=0:5,samples=120] plot ({\x},{2.5*exp(-0.45*\x)});
\draw[domain=0:5,samples=120] plot ({\x},{2.2*exp(-0.5*\x)});
\draw[dashed,domain=0:5,samples=120] plot ({\x},{2.4*exp(-0.8*\x)});
\draw[dotted,domain=0:5,samples=120] plot ({\x},{0.9*exp(-0.3*\x)});
\node[right] at (2.3,1.15) {\footnotesize \(g\)};
\node[right] at (3.3,0.32) {\footnotesize \(f_n\)};
\node at (2.6,-0.9) {\footnotesize 有可积包络：面积无处可逃};
\begin{scope}[shift={(8.4,0)}]
\draw[->] (0,0) -- (5.9,0) node[right] {\footnotesize \(x\)};
\draw[->] (0,0) -- (0,3.5);
\draw[thick,dotted] (0,0.8) -- (5,0.8) -- (5,0);
\draw[thick,dashed] (0,1.6) -- (2.5,1.6) -- (2.5,0);
\draw[thick] (0,3.2) -- (1.25,3.2) -- (1.25,0);
\begin{scope}
\clip (0,0) rectangle (5.9,3.4);
\draw[very thick,domain=0.29:5,samples=140] plot ({\x},{1.0/\x});
\end{scope}
\node[right] at (2.0,0.72) {\footnotesize \(\sup_n f_n\sim 1/x\)};
\node at (2.6,-0.9) {\footnotesize 包络本身不可积：拦不住};
\end{scope}
\end{tikzpicture}
\end{center}

\emph{一片固定的可积天花板压住整族函数，尖峰就长不起来，质量也漂不远。}

\subsection{一个与 \(n\) 无关的可积上界就够了}

\begin{theorem}[控制收敛定理]
设 \(f_n\) 可测且几乎处处有 \(f_n(x)\to f(x)\)。若存在\emph{与 \(n\) 无关}的可积函数 \(g\) 使
\[
\abs{f_n(x)}\le g(x)\qquad\text{对每个 }n\text{ 几乎处处成立},
\]
则 \(f\) 可积，且 \(\int\abs{f_n-f}\to0\)，特别地 \(\int f_n\to\int f\)。
\end{theorem}

结论比``两个数值相等''强。\(\int\abs{f_n-f}\to0\) 说的是 \(L^1\) 收敛，被积函数本身在积分意义下靠拢，而不只是两侧积分值凑巧对上。这条更强的结论在后续估计中经常直接被引用。

左图里的三条 \(f_n\) 各有各的形状，可以上下起伏、可以正负交替，但没有一条能钻出 \(g\) 划定的范围。\(g\) 的积分有限，意味着整个可用面积从一开始就是有限的且分布固定——上一节点出的两条逃逸通道，高度无界飙升与支撑集漂向远方，被同一个上界一起封死。右图的对照说明这个上界必须自己可积：尖峰族的逐点上确界确实存在，只是它的积分是无穷，等于没设上限。

顺带说一句，共同上界是常数也不一定安全。积分域无限时，常数函数的积分发散，同样不是合格的包络。``可积''管的是总量，不是局部大小。

\subsection{三项条件各防住一种失败}

条件对账值得写全，因为实际使用中出错几乎都发生在漏掉某一项。

可测性保证 \(\int f_n\) 和 \(\int f\) 这些记号有意义。它看着像技术性前提，作用是把 \(f\) 关进可以谈论积分的函数类；逐点极限自动保持可测，所以这一项通常在写下 \(f_n\) 时就已经满足，但换成不可测的构造时结论会直接失去主语。

逐点收敛提供极限函数本身。没有它，你连要与什么比较都不知道；有它但只有它，就退回上一节的处境——尖峰列的逐点收敛毫无问题，塌的是别处。

可积控制函数是唯一带来新信息的那一项，它的两个修饰词缺一不可。``与 \(n\) 无关''排除了随序列一起变大的伪上界；``可积''排除了总量无穷的空头承诺。有了它，\(\abs{f_n-f}\le2g\) 成立，Fatou 引理用在非负函数 \(2g-\abs{f_n-f}\) 上即得 \(\int\abs{f_n-f}\to0\)，证明骨架就这三行——可以看出定理的全部力气都花在把非单调的情形化归成可以用 Fatou 的情形。

\subsection{在概率空间里找包络要容易得多}

期望是对概率测度的积分，全空间的测度等于 1。这个前提让控制条件的寻找简单了一大截：常数函数在概率空间上可积，所以只要 \(\abs{X_n}\le M\) 且 \(X_n\to X\) 几乎处处，控制收敛定理立刻给出 \(\E[X_n]\to\E[X]\)。更一般地，\(\abs{X_n}\le Y\) 且 \(\E[Y]<\infty\) 同样够用。参数估计的大样本性质、渐近方差的推导、以及蒙特卡洛估计中``先取极限还是先取期望''的合法性，背后都是这一步。

条件被突破时的失效方式也很典型。取
\[
X_n=\begin{cases}
n,&\text{以概率 }1/n,\\
0,&\text{否则},
\end{cases}
\]
则 \(X_n\to0\) 依概率成立，而 \(\E X_n=1\) 不动。小概率的大值扛住了全部期望质量，任何可积的 \(Y\) 都压不住这一列——这与上一节的尖峰是同一张图，只是横轴换成了样本空间。训练中偶发的极端损失值让平均损失长期偏离真实期望，属于同一类现象。

把包络条件放宽的标准做法是改用一致可积性，配上依概率收敛即可交换极限与期望，这是 Vitali 收敛定理的内容，\hyperref[sec:12-Real-Analysis-Support/05-Measure-Integration-and-Conditional-Expectation/06]{``一致可积性阻止期望质量逃逸''}把这一条件展开来讲。

\subsection{同一套条件平移到无穷级数上}

无穷级数是部分和的极限，\(\sum_{k\ge1}f_k=\lim_n\sum_{k\le n}f_k\)，所以求和号与积分号能否交换属于同一个问题。若 \(f_k\ge0\)，部分和单调上升，单调收敛定理直接给出
\[
\int\sum_{k=1}^{\infty}f_k=\sum_{k=1}^{\infty}\int f_k,
\]
两侧允许为无穷，这是 Tonelli 定理的核心情形。

放弃非负性后，替代条件是绝对可积：\(\sum_k\int\abs{f_k}<\infty\)。此时部分和被 \(\sum_k\abs{f_k}\) 控制住，用控制收敛定理或 Fubini 都能得到交换且结果有限。只有条件收敛时最危险，重排项就可能改变答案，所以每次把积分号推进求和号之前，都该说清楚依据的是非负性还是绝对可积性。

\subsection{它是充分条件，不是判据}

存在积分确实收敛却找不到简洁可积包络的情形，那时一致可积性或针对具体参数族的估计是更精细的替代路径。控制收敛定理的价值在于条件好检验、适用面宽，它从不宣称刻画了所有可交换的情况。

使用中最常见的无效操作是取 \(g_n=\abs{f_n}\) 充当控制函数。这个式子恒成立，也恒等于什么都没说——它依赖 \(n\)，没有提供任何超出原序列的信息。包络必须是一个固定的、在极限过程开始之前就摆在那里的函数。

到此，逐点极限穿过积分号有了三条通道：一致收敛让误差整体消失，单调收敛禁止质量迁移，控制收敛给整族加一个可积上界。三条通道守的是同一条底线——总质量不许从任何缝隙溜走。下一节把同样的思路用到差商上，那里要控制的对象换了，逻辑一模一样。
